\section{What This Manual Assembles}

The 1962 Roman Missal is not assembled by opening to a saint, copying ten variable texts, and inserting them wherever a hand missal leaves a blank. Before a formulary can be joined to the stable Order of Mass, the calendar must identify the liturgical day, the table of precedence must resolve any occurrence, the rules must determine what happens to the impeded item, and the Missal must permit the proposed kind of Mass. Only then can the proper, Common, seasonal interventions, commemorations, and ritual directions be placed into the \latin{Ordo Missae}.

This manual explains that process for the Latin \latin{Missale Romanum}, editio typica, 1962. Its controlling law is the 1960 \latin{Codex rubricarum}, promulgated by the Sacred Congregation of Rites after Saint John XXIII approved it in \latin{Rubricarum instructum}, and incorporated into the 1962 book. References simply to a numbered rubric mean the continuous numbering of that Code. References to the \latin{Ritus servandus}, \latin{Ordo Missae}, \latin{Canon Missae}, or a particular formulary name the other controlling part of the Missal.

\begin{studybox}{The five questions, in their necessary order}
\begin{enumerate}[leftmargin=1.45em]
\item \textbf{What day occurs?} Build the temporal, universal sanctoral, and applicable particular calendars.
\item \textbf{Which day has precedence?} Locate each competing day in rubric 91's twenty-eight-position table and apply the express Sunday and octave exceptions.
\item \textbf{What becomes of the impeded item?} Omission, commemoration, accidental transfer, and perpetual reposition are different legal results.
\item \textbf{What kind of Mass is admitted?} The Mass corresponding to the Office, a festive Mass, a votive Mass, a Requiem, a ritual Mass, and an external solemnity obey different rules.
\item \textbf{How is the admitted Mass assembled?} Insert its proper texts, allowed commemorations, seasonal changes, and special directions into the stable Order.
\end{enumerate}
\end{studybox}

\subsection{Object and reader}

The immediate reader is a priest, server, sacristan, schola director, editor, or serious lay student who needs to know where a day's texts come from and why. This is a textual and rubrical preparation manual, not a complete ceremonial. It distinguishes Low Mass, sung Mass, and Solemn Mass when the distinction changes an admitted commemoration, tone, minister, reading, or audible part; it does not replace the \latin{Ritus servandus}, a competent ceremonial, or the direction of ecclesiastical authority.

The word \emph{proper} is used in two related senses. A proper calendar belongs to a diocese, place, church, or religious institute. A proper text is variable material appointed for a celebration. Context will identify which is meant. The \emph{Ordinary} here means the stable textual and ritual framework of the complete Order of Mass, not merely Kyrie, Gloria, Credo, Sanctus, and Agnus Dei.

\subsection{Edition boundary}

This manual deliberately excludes both earlier and later systems. It does not restore the pre-1960 categories of double, semidouble, simple, privileged octave, transferred Sunday, or multiplied commemorations. It does not import the postconciliar ranks of solemnity, feast, memorial, optional memorial, or the rules of a later General Instruction. A familiar practice is not evidence that the 1962 book commands it.

Present canonical authorization to celebrate using the 1962 Missal is mutable and distinct from the historical book's internal assembly rules. This work makes no claim that a person, place, date, or proposed celebration is presently authorized. It also cannot supply an unknown diocesan calendar, church title, anniversary of dedication, patron, religious calendar, or indult. A competent Ordo remains indispensable precisely because it applies the universal rules to those local facts.

\subsection{Method and authority}

The official Latin rules control. The English explanations are original study paraphrases, not an approved translation. Exact rubric numbers are provided so that the reader can return to the official text. The published 1962 Missal facsimile controls sequence and formulary directions; searchable OCR and calendar software may locate a passage or test a calculation, but they cannot override the page image or create a local law that has not been sourced.

Three labels govern the exposition:

\begin{itemize}
\item A \textbf{direct rule} restates an express numbered rubric or formulary direction.
\item A \textbf{necessary application} combines express rules to answer a case, without creating a new permission.
\item An \textbf{editorial workflow} arranges the rules into a usable checklist; its order is pedagogical, not another promulgated rubric.
\end{itemize}

The worked cases use conditions rather than invented civil dates unless the date is inherent in the rule. A hypothetical “parish titular” assumes that the title and rank are established in a competent proper calendar; the example does not establish them for an actual church.

\subsection{The governing principle}

Rubric 270 states the ordinary relation: Mass should agree with the Office of the day. The exceptions do not abolish that relation. They are enumerated paths---festive, votive, Requiem, ritual, or externally solemn---whose admission must be proved before their texts are assembled. The safest question is therefore not “Which Mass would be fitting?” but “Which source makes this Mass available on this liturgical day, and what exactly does that source allow?”
